\begin{casebox}
\textbf{知识点小案例：局部 top-k 不能推出全局 top-k。}
把输入切成 10 个 chunk，每个 chunk 只保留局部 top 3。构造一个 key：它在每个 chunk 都排第 4，但全局累加后超过所有局部 top key。若 reduce 阶段只看局部 top 3，这个 key 永远消失。由此推出规律：内核模式必须说明是否丢信息，不能只说明“减少数据量”。工程上 top-k 可以是精确算法，也可以是近似算法，但必须写明候选覆盖规则、误差合同和反例输入；否则一个很快的局部优化会破坏全局语义。
\end{casebox}
